%% kapa: exam=1 round=31 year=2020 period=2 form=A subject=accounting subject_ko=회계학 seq=29 number=69
%% source: Q-Net 감정평가사 기출문제 (cst00309) artlSeq=5207219 file=extracted/1008/5207219/A형/A형/2020년도 제31회 감정평가사 1차 시험_2교시_A형.pdf
%% conversion: kordoc 4.24.51 -> tools/gen1.py (auto); review=proofread-sonnet; flags=none
\begin{question}{69}
상품매매기업인 (주)감평은 20x0년 말 취득한 건물(취득원가 ￦2,400,000, 내용연수 10년, 잔존가치 ￦0, 정액법 상각)을 유형자산으로 분류하여 즉시 사용개시 하고, 동 건물에 대해 재평가모형을 적용하기로 하였다. 20x1년 10월 1일 (주)감평은 동 건물을 투자부동산으로 계정 대체하고 공정가치모형을 적용하기로 하였다. 시점별 건물의 공정가치는 다음과 같다.
\par{\small\centering\begin{tabular}{|c|c|c|}\hline
20x0년 말 & 20x1년 10월 1일 & 20x1년 말 \\ \hline
￦2,400,000 & ￦2,300,000 & ￦2,050,000 \\ \hline
\end{tabular}\par}
동 건물 관련 회계처리가 20x1년 당기순이익과 기타포괄이익에 미치는 영향은 각각 얼마인가? (단, 재평가잉여금은 이익잉여금으로 대체하지 않으며, 감가상각은 월할계산한다.)
{\par\noindent\hspace*{2em}\makebox[4cm][c]{\uline{당기순이익}}\makebox[4cm][c]{\uline{기타포괄이익}}\par}
\begin{choices}
  \item \makebox[4cm][c]{￦180,000 감소}\makebox[4cm][c]{￦80,000 증가}
  \item \makebox[4cm][c]{￦180,000 감소}\makebox[4cm][c]{￦350,000 증가}
  \item \makebox[4cm][c]{￦430,000 감소}\makebox[4cm][c]{￦80,000 증가}
  \item \makebox[4cm][c]{￦430,000 감소}\makebox[4cm][c]{￦350,000 증가}
  \item \makebox[4cm][c]{￦430,000 감소}\makebox[4cm][c]{￦430,000 감소}
\end{choices}
\end{question}
\answer{3}
